% LaTeX file for resume 
% This file uses the resume document class (res.cls)
\RequirePackage{cmap}
\documentclass{res} 
%\usepackage{helvetica} % uses helvetica postscript font (download helvetica.sty)
%\usepackage{newcent}   % uses new century schoolbook postscript font 
\setlength{\textheight}{9.5in} % increase text height to fit on 1-page 

\usepackage{hyperref}
\usepackage[T1]{fontenc}
\usepackage{lmodern}
\usepackage{cmap}
\usepackage[utf8]{inputenc}
\usepackage[english,russian]{babel}

\begin{document} 

\name{Sophia Kovaleva\\[12pt]}     % the \\[12pt] adds a blank
				        % line after name      

\address{Phone: (302) 415-4176\\Email, Google Hangouts: \href{mailto:sophia.m.kovaleva@gmail.com}{sophia.m.kovaleva@gmail.com}}


\begin{resume}

\section{SUMMARY}

I'm a sixth year PhD student in Electrical and Computer Engineering at Carnegie Mellon University working on transparency and fairness in machine learning and expecting to graduate in April -- May 2019.
I'm looking to full-time employment opportunities in the area of AI and ML research (particularly, recommender systems and neural networks), data science and analytics, privacy engineering, and software engineering. 
My core skills include designing and implementing machine learning systems that are subject to specified fairness and privacy constraints and are working with user data.


\section{EDUCATION}
\begin{itemize}
	\item
		\underline{Carnegie Mellon University Silicon Valley Campus, July 2013 -- May 2019 (expected)}\\
		PhD student, Electrical and Computer Engineering. Advisor: Dr. Anupam Datta, PhD thesis: ``Representation-aware accounting of machine learning models''.
	\item
		\underline{Carnegie Mellon University Silicon Valley Campus, July 2013 -- February 2019}\\
		Master of Science in Electrical and Computer Engineering
	\item
		Moscow State Institute of Electronics and Mathematics (MIEM), September 2007 -- June 2012\\
		Engineer degree in Computational Machines, Complexes, Systems and networks\\
		G.P.A. 4.9 out of 5.0
	\item
		English Language Institute of the University of Delaware, July 2012 -- August 2012\\
		English language practice, Certificate at the High Advanced Level with Honors
\end{itemize}

\section{SKILLS}          
\begin{itemize}
	\item
		Programming in Python, Java (Android), some JavaScript and C/C++.
	\item
		Building machine learning systems with TensorFlow, Apache Spark, scikit-learn, NumPy.
	\item
		Data visualization with matplotlib, GNU Octave, D3.js + Django REST.
	\item
		General GNU/Linux (especially Ubuntu and Gentoo) usage and administration including networking and virtualization.
	\item
		Experience in GNU Octave, Maple, Haskell, Pascal, Fortran, shell scripting.
	\item
		Experience with online video broadcasting/conferencing and video production.
\end{itemize}
 
\section{EXPERIENCE}
\vspace{-0.1in}	
\begin{tabbing}
\hspace{2in}\= \hspace{2in}\= \kill % set up two tab positions
{\bf PhD Student / RA} \>Carnegie Mellon University \>Fall 2013 -- Present\\
			\>Moffett Field, CA \> 
\end{tabbing}      % suppress blank line after tabbing
Notable research projects:
\begin{itemize}
	\item
		Room Occupancy Detection in SmartSpace -- using data from wall-mounted physical sensors to estimate the presence and the number of people in rooms.
	\item
		Indoor Positioning of Mobile Devices by Combined Wi-Fi and GPS Signals -- improving fingerprinting techniques for indoor positioning using continuous Bayesian classifier.
	\item
		Improving Wireless Emergency Alerts -- implementing location-based client-side filtering of broadcast emergency messages to mobile phones to reduce the number of irrelevant messages seen by users, and improve user experience.
	\item
		Wireless communication set-up for RSSI broadcast -- building 3D radiation pattens of cellular antennae by interpolating RSSI data collected by a drone-based platform.
	\item
		Explaining matrix factorization recommender systems with supervised learning.
	\item
		Removing the use of sensitive attributes (race, gender, etc.) from the decision process of machine learning systems.	
\end{itemize}

\vspace{-0.1in}	
\begin{tabbing}
\hspace{2in}\= \hspace{2in}\= \kill % set up two tab positions
{\bf Data Scienece Intern} \>Zugata, Inc. \>Summer 2015\\
			\>Palo Alto, CA \> 
\end{tabbing}      % suppress blank line after tabbing
Building and analyzing corporate social graphs based on metadata about users' communications. My job involved full stack development with Python/Django in the backend and JS/d3 in the frontend.

\vspace{-0.1in}	
\begin{tabbing}
\hspace{2in}\= \hspace{2in}\= \kill % set up two tab positions
{\bf Technician / TA} \>MIEM / MIEM HSE \>Fall 2010 -- June 2013\\
			\>Moscow, Russia \> (volunteered from 2008)
\end{tabbing}      % suppress blank line after tabbing
Summary for multiple years:
\begin{itemize}
	\item
		Arranging materials and teaching some of the sessions of a programming class to freshmen. Languages: Python, C. Topics: syntax, basic algorithms with lists and strings (Python), developing 2D games using the PyGame library, basic algorithms with arrays, strings and matrices (C), dynamic memory allocation, linked lists, trees and graphs.
	\item
		Developing and teaching an applied programming class to sophomores. Language: Python. Topics: Python object model, basic functional programming, multiprocessing, socket network interface, GTK-based GUIs, databases and others.
\end{itemize}

\section{LANGUAGES SPOKEN}
\begin{enumerate}
	\item
		Russian (native);
	\item
		English (114 TOEFL iBT);
	\item
		Japanese (JLPT level N4);
	\item
		French (basic).
\end{enumerate}
 
\section{PUBLICATIONS}
\begin{itemize}
	\item
		Datta, A. et al (2018)  Latent Factor Interpretations for Collaborative Filtering, \url{https://arxiv.org/abs/1711.10816}.
	\item
		Erdogmus, H. et al (2015) Opportunities, Options, and Enhancements for the Wireless Emergency Alerting Service, Department of Homeland Security, First Responders Group, February 2016.  
	\item
		Kovalev, M. (2014) Indoor Positioning of Mobile Devices by Combined Wi-Fi and GPS Signals. \emph{Indoor Positioning and Indoor Navigation (IPIN), 2014 International Conference on}, pp.332-339, 27-30 Oct. 2014, doi: 10.1109/IPIN.2014.7275500
	\item
		Kovalev, M. (2013). Testing and the analysis of the perforamce of an algorithm for positioning public transportation vehicles by GSM signal. \emph{Proceedings of the annual scientific and technical conference of undergraduate and graduate students, and young experts of the Moscow State Institute of Electronics and Mathematics, National Research University Higher School of Economics} (pp. 140--141). Moscow, Russia: MIEM NRU HRE Press. ISBN 978-5-94768-066-9.\\
		(М.М.Ковалев. Испытания и анализ производительности алгоритма позиционирования наземного транспорта по сигналу GSM. Научно-техническая конференция студентов, аспирантов и молодых специалистов МИЭМ НИУ ВШЭ. Тезисы докладов. -- М.: МИЭМ НИУ ВШЭ, 2013. -- 316 с. ISBN 978-5-94768-066-9. с.с. 140--141)
	\item
		Kovalev, M. (2012). Prototype of the software and hardware complex for collecting data about GSM signal levels. \emph{Proceedings of the annual scientific and technical conference of undergraduate and graduate students, and young experts of the Moscow State Institute of Electronics and Mathematics} (pp. 185--187). Moscow, Russia: MIEM Press. ISBN 978-5-94506-314-3.\\
		(М.М.Ковалев. Прототип программно-аппаратного комплекса для сбора данных об уровнях сигнала GSM на местности. Научно-техническая конференция студентов, аспирантов и молодых специалистов МИЭМ, посвящённая 50-летию МИЭМ. Тезисы докладов. -- М.: МИЭМ, 2012. --- 445 с. ISBN 978-5-94506-314-3. с.с. 185--187.)
	\item
		Kovalev, M. (2011). Positioning of the surface transport vehicles by means of GSM signal. \emph{Proceedings of the annual scientific and technical conference of undergraduate and graduate students, and young experts of the Moscow State Institute of Electronics and Mathematics} (pp. 29--30). Moscow, Russia: MIEM Press. ISBN 978-5-94506-257-3.\\
		(М.М.Ковалев. Позиционирование наземного транспорта с помощью сигнала GSM. Научно-техническая конференция студентов, аспирантов и молодых специалистов МИЭМ. Тезисы докладов. -- М.: МИЭМ, 2011. --- 420. ISBN 978-5-94506-257-3. с.с. 29--30.)
	\item
		Kovalev, M. (2008). Cyclic development of Russia \emph{Proceedings of the annual scientific and technical conference of undergraduate and graduate students, and young experts of the Moscow State Institute of Electronics and Mathematics} (pp. 398--399). Moscow, Russia: MIEM Press. ISBN 978-5-94506-170-5.\\
		(М.М.Ковалев. Циклическое развитие России. Научно-техническая конференция студентов, аспирантов и молодых специалистов МИЭМ. Тезисы докладов. -- М.: МИЭМ, 2008. --- 495. ISBN 978-5-94506-170-5. с.с. 398--399.)
\end{itemize}

\section{NOTABLE COURSEWORK}          
 
\begin{itemize}
	\item
		Artificial intelligence (online by Peter Norvig and Sebastian Thrun)
	\item
		Machine learning (online by Andrew Ng)
	\item
		Special course: fundamentals of neuroscience (public class, Moscow State University)
	\item
		Statistical Discovery and Learning (18-799, Carnegie Mellon University)
	\item
		Machine Learning for Signal Processing (18-797, Carnegie Mellon University)
	\item
		Special Topics in Signal Processing: Advanced Machine Learning (18-799 SM, Carnegie Mellon University)
	\item
		Algorithms: Design and Analysis, Part 1 (online, Coursera)
	\item
		Exploring Neural Data (online, Coursera)
	\item
		Model Thinking (online, Coursera)
	\item
		Probability theory, mathematical statistics and stochastic processes  (MIEM)
	\item
		7.00x Intro to Biology -- The Secret of Life (online, edX)
\end{itemize}

\section{MISCELLANEOUS}

Received a grant for my undergraduate thesis research in 2011.\\
Amateur radio operator KK6HAI.

More detailed information about my coursework, skills, experience, education, publications, research projects, honors, and awards is available in my LinkedIn profile:\\ \href{https://www.linkedin.com/in/sophia-kovaleva-1a664a3b/}{https://www.linkedin.com/in/sophia-kovaleva-1a664a3b/}.

My portfolio can be found on my GitHub page: \href{https://github.com/kosonya/}{https://github.com/kosonya/}.

\end{resume}
\end{document}
